\section*{Physique (13 points)}
\section*{Exercice 1(4 points) : propagation d'une onde}
Durant des séances de travaux pratiques, des élèves ont procédé à :

\begin{itemize}
  \item l'étude de la propagation d'une onde mécanique progressive périodique à la surface de l'eau ;
  \item la détermination de la vitesse de propagation du son dans la salle de TP ;
  \item la détermination de la longueur d'onde d'une onde lumineuse monochromatique.
\end{itemize}

\section*{1. Propagation d'une onde à la surface de l'eau}
On produit à l'aide d'une plaque $(P)$ d'un vibreur, à la surface libre de l'eau d'une cuve à ondes, des ondes progressives périodiques de fréquence $N=10 \mathrm{~Hz}$. Les ondes se propagent sans amortissement ni réflexion. La figure (1) donne l'aspect de la surface de l'eau à un instant donné.\\
Donnée : $d=6 \mathrm{~cm}$.\\
0,5 1.1. Déterminer la valeur de la longueur d'onde $\lambda$.\\
0,5\\
1.2. Déduire la valeur de la vitesse de propagation $v$ à la surface de l'eau.\\
\includegraphics[width=\textwidth]{2025_03_20_cb61d7da9b7d1381f8fag-3(1)}

Figure 1

0,5 1.3. On considère deux points $M$ et $P$ de la surface de l'eau, tel que $M P=7 \mathrm{~cm}$ (figure 1).\\
Calculer le retard temporel $\tau$ de la vibration du point $P$ par rapport à $M$.

\section*{2. Détermination expérimentale de la vitesse de propagation du son}
Pour déterminer la vitesse de propagation d'une onde sonore dans la salle de TP, l'enseignant a préparé le montage expérimental de la figure (2) (page 4/6) qui comporte :

\begin{itemize}
  \item deux microphones $M_{1}$ et $M_{2}$ séparés par une distance $d$;
  \item un oscilloscope ;
  \item un haut-parleur ;
  \item un GBF réglé à une fréquence $N$.
\end{itemize}

La figure (3) (page 4/6) donne les oscillogrammes observés pour une distance $d_{1}=21 \mathrm{~cm}$.\\
La sensibilité horizontale est $S_{h}=1,0.10^{-4}{\mathrm{~s} . \mathrm{div}^{-1}}$.



\begin{center}
\includegraphics[width=\textwidth]{2025_03_20_cb61d7da9b7d1381f8fag-4(1)}
\end{center}

0,5 2.1. Déterminer la valeur de la période $T$ de l'onde sonore.\\
2.2. On déplace horizontalement le microphone $M_{2}$ progressivement par rapport à $M_{1}$ jusqu'à ce que les deux courbes soient à nouveau en phase. La distance entre les deux microphones est alors $d_{2}=41,5 \mathrm{~cm}$.\\
0,5 a. Déterminer la valeur de la longueur d'onde $\lambda$ de l'onde sonore.\\
$\mathbf{0 , 5}$ b. Calculer la valeur de la vitesse de propagation $v$ du son dans l'air.\\
3. Détermination de la longueur d'onde d'une onde lumineuse

Pour déterminer la longueur d'onde d'une onde lumineuse, les élèves ont éclairé une fente de largueur $a=5,0.10^{-5} \mathrm{~m}$ par un faisceau de lumière monochromatique. Ils ont observé des taches lumineuses sur un écran situé à la distance $D=1,5 \mathrm{~m}$ de la fente (Figure 4). La mesure de la largeur de la tache centrale a donné $L=3,8 \mathrm{~cm}$.\\
3.1. Nommer le phénomène observé durant cette expérience.

0,75\\
3.2. Établir l'expression de la longueur d'onde $\lambda$ en fonction de $L, D$ et $a$ (On considère que $\tan \theta \approx \theta(\mathrm{rad})$ ). Calculer $\lambda$.\\
\includegraphics[width=\textwidth]{2025_03_20_cb61d7da9b7d1381f8fag-4}

Figure 4

\section*{Exercice 2 ( 2,5 points) : Le radon et la qualité de l'air}
La Terre émet de façon naturelle le gaz radon. Ce gaz qui se propage facilement à l'intérieur des immeubles est radioactif. Il est considéré comme l'une des principales causes du cancer $\alpha$ du poumon après la cigarette. Selon l'instance internationale de la protection radioactive, la concentration volumique de la radioactivité du gaz radon dans l'air des locaux ne doit pas dépasser $400 \mathrm{~Bq} . \mathrm{m}^{-3}$.

\section*{Données :}
\begin{center}
\begin{tabular}{|l|c|c|c|c|}
\hline
Noyau & Francium & Radon & Polonium & Hélium \\
\hline
Symbole & ${ }_{87}^{223} \mathrm{Fr}$ & ${ }_{86}^{222} \mathrm{Rn}$ & ${ }_{84}^{28} \mathrm{Po}$ & ${ }_{2}^{4} \mathrm{He}$ \\
\hline
Masse du noyau en unité (u) & 222,9720 & 221,9704 & 217,9628 & 4,0015 \\
\hline
\multicolumn{4}{|c|}{$1 u=931,5 \mathrm{MeV} . \mathrm{c}^{-2}$} &  \\
\hline
\end{tabular}
\end{center}

\begin{enumerate}
  \item Donner la composition du noyau de radon ${ }_{86}^{222} R n$.
  \item Écrire l'équation de désintégration du radon ${ }_{86}^{222} R n$, en précisant le noyau fils.
  \item Calculer en unité $(\mathrm{MeV})$, la valeur de l'énergie libérée $E_{\text {liberree }}=|\Delta E|$ au cours de la désintégration d'un noyau de radon ${ }_{86}^{222} \mathrm{Rn}$.\\
\includegraphics[width=\textwidth]{2025_03_20_cb61d7da9b7d1381f8fag-5(1)}
\end{enumerate}

Les bobines et les condensateurs sont très utilisés dans les circuits des appareils électriques et électroniques répandus comme les jouets des enfants, les montres électriques, les systèmes d'alarme et appareils de contrôle...L'analyse de ces circuits peut se faire par une étude électrique ou énergétique, ce qui permet de déterminer certaines grandeurs caractéristiques et mettre en évidence les échanges énergétiques qui se produisent.\\
Cet exercice vise :

\begin{itemize}
  \item la détermination des grandeurs ( $L ; r$ ) caractérisant une bobine ;
  \item l'étude des oscillations électriques libres dans un circuit $R L C$ série.
\end{itemize}

\section*{Partie 1 : détermination des grandeurs ( $L ; r$ ) caractérisant une bobine}
Un professeur met à la disposition des élèves les outils suivants :

\begin{itemize}
  \item une bobine $r$ et de résistance $L$ d'inductance (b) ;
  \item un condensateur de capacité $C$;
  \item un conducteur ohmique de résistance $R=90 \Omega$;
  \item un générateur $G_{1}$ de force électromotrice $E=6 \mathrm{~V}$;
  \item un générateur idéal de courant $G_{2}$;
  \item un interrupteur $K$;
  \item un oscilloscope ;
  \item des fils de connexion.
\end{itemize}

\begin{enumerate}
  \item Citer parmi les outils signalés précédemment, ceux nécessaires à la réalisation d'un circuit permettant d'étudier la réponse d'un dipôle $R L$ à un échelon de tension ascendant. 2. Quel est le rôle de la bobine lors de la fermeture du circuit?
  \item Établir l'équation différentielle vérifiée par l'intensité du courant $i(t)$ traversant le circuit.
  \item Sachant que la solution de l'équation différentielle s'écrit: $i(t)=\mathrm{I}_{0}\left(1-\mathrm{e}^{\frac{-t}{\tau}}\right)$, déterminer les expressions de $\mathrm{I}_{0}$ et $\tau$ en fonction des paramètres du circuit.
  \item À l'aide d'un système d'acquisition informatique, les élèves ont obtenu la courbe représentée sur la figure (1) .\\
\includegraphics[width=\textwidth]{2025_03_20_cb61d7da9b7d1381f8fag-5}
\end{enumerate}

Figure 1


\includegraphics[width=\textwidth]{2025_03_20_cb61d7da9b7d1381f8fag-6}

0,5 a. Déterminer graphiquement les valeurs de $\mathrm{I}_{0}$ et de $\tau$.\\
0,5\\
b. Vérifier que $r=10 \Omega$ et $L=1 H$.\\
c. Déterminer la valeur de la tension $u_{b}$ aux bornes de la bobine en régime permanent.

\section*{Partie 2: Oscillations électriques libres dans un circuit $R L C$ série}
Après avoir chargé totalement le condensateur cité dans la liste des outils, les élèves procèdent à sa décharge à travers la bobine (b). La courbe de la figure (2) représente les variations de la tension $u_{C}(t)$ aux bornes du condensateur au cours de sa décharge.\\
\includegraphics[width=\textwidth]{2025_03_20_cb61d7da9b7d1381f8fag-6(1)}

Figure 2\\
0,5 1. Représenter le schéma du montage expérimental permettant de réaliser la décharge du condensateur.\\
2. Déterminer graphiquement la valeur de la pseudo- période $T$.

En déduire la valeur de $C$. On considère que la valeur de la pseudo- période $T$ est égale à la période propre. $\pi^{2}=10 \mathrm{n}$ prend. $\mathrm{O}(L C)$ de l'oscillateur $T_{0}$\\
3. Interpréter l'allure de la courbe de point de vue énergétique.\\
4. Sous quelle forme est emmagasinée l'énergie dans le circuit à l'instant $t=\frac{T}{4}$ ? Justifier votre réponse.\\
5. Calculer la variation de l'énergie totale $\Delta \mathcal{E}$ du circuit entre les instants $t_{0}=0$ et $t_{1}=4 T$.\\
6. Pour entretenir les oscillations électriques, on ajoute dans le circuit ( $R L C$ ), un générateur $G$ qui donne une tension $u_{G}$ proportionnelle à l'intensité du courant qui traverse le circuit ( $u_{G}=k . i$ ).\\
a. Indiquer le rôle du générateur $G$ de point de vue énergétique.

0,25\\
b. Déterminer la valeur de $k$ pour que le circuit soit le siège d'oscillations électriques entretenues.



